%
% Dokumentation
%

% Allgemeine Dokumenteigenschaften
\documentclass[a4paper,			% Druck auf DIN A4-Papier.
numbers=enddot,					% Die Kapitel bekommen einen Punkt.
chapterprefix=false,			% Das Wort "Kapitel:" wird gelöscht.
12pt,							% 12pt ist die Normalschriftgröße, weitere Möglichkeitenn sind 10pt oder 11pt.
pointlessnumbers				% Nummerierungen in den Kapiteln haben keinen finalen Punkt zb 1.1 statt 1.1.
]							
{scrreprt}  					% Dokumentart.

\setkomafont{chapter}{\Large}
\setkomafont{section}{\large}

% Präambel: Parameter die sich auf das Gesamte Dokument auswirken.
\usepackage{geometry}
\usepackage{float}
\usepackage[ngerman]{babel}				% "Neue" Deutsche Rechtschreibung
\usepackage{textcomp}					% Aktiviert das Eurosymbol. (Im Text schreiben mit \texteuro)
\usepackage[utf8]{inputenc}				% Direkte Eingabe von Umlauten im Text. (Dieser Code läuft nicht überall)
\usepackage{graphicx}					% Ermöglicht das anzeigen von Bildern
%\usepackage{babelbib}					% Für verschiedene Deutsche Stile des Literaturverzeichnis
\usepackage{apacite}
%\usepackage[babel, german=guillemets]{csquotes}
\usepackage[babel,german=quotes]{csquotes}
\usepackage{amsmath,amssymb}			% Das Mathe Paket "American Mathematic Society" laden
\usepackage[printonlyused]{acronym}		% Paket für das Abkürzungsverzeichnis
\usepackage{fancyhdr}					% Fancyhdr zum Nutzen von Kopf und Fußzeilen mit voreingestellten Formatierungen
\usepackage{xcolor}						% Verschiedene Farben sind direkt als white usw nutzbar
\usepackage{epstopdf}					% eps Postscript Dateien können direkt eingefügt werden
\usepackage{tocstyle}					% Das Inhaltsverzeichnis lässt sich konfigurieren
\usepackage{listings}					% Ermöglicht das Einfügen von Quellcode mit Kommentarhighlighting und Nummerierung aussen
\usepackage{subfig}						% 2 Bilder nebeneinander
\usepackage{chngcntr}
\usepackage{tabularx}
\counterwithout{figure}{chapter} 




% Kapitel, Bilder, Abkürzungen usw werden zu Links, ohne roten Rahmen. DAS FUNKTIONIERT NICHT MIT APACITE!!!!
%\usepackage[plainpages=false,
%pageanchor=false,
%linkbordercolor=white,					% Die Rahmen werden weiß und sind somit nicht mehr im PDF sichtbar
%pdfpagelabels]
%{hyperref}								% Inhaltsverzeichnis wird als Lesezeichen im Adobe Reader angezeigt	

% Absatz-Einstellungen
\geometry{a4paper,
left=20mm,
right=20mm,
top=2.5cm,
bottom=2.5cm}
\parindent0ex							% Oberster Satz eines Absatzes wird NICHT eingerückt!
\parskip1.4ex plus0.2ex minus0.2ex


% Kopf und Fußzeile einbinden und festlegen.
\pagestyle{fancy}									% Eigenen Seitenstil verwenden
\fancyhf{} 											% clear all header and footer fields
\renewcommand{\headrulewidth}{1pt}					% Die Dicke der Kopfzeilen-Trennline wird auf 1pt festgelegt
\renewcommand{\footrulewidth}{1pt}					% Die Dicke der Fusszeilen-Trennline wird auf 1pt festgelegt
\fancypagestyle{plain}{}
\fancyhead[R]{}										% Kopfzeile Rechts
\fancyfoot[R]{\thepage}								% Fusszeile Rechts

% Inhaltsverzeichnis anpassen.
\usetocstyle{allwithdot}				% Fällt den Leerraum mit Punkten auf
\setcounter{tocdepth}{1}

% Das eigentliche Dokument.
\begin{document}

% Dokumentenweite Einstellungen und Kommandos
\nonfrenchspacing
\newcommand{\FigRef}[1]{Abbildung \ref{#1} auf Seite \pageref{#1}}
\newcommand{\TabRef}[1]{Tabelle \ref{#1} auf Seite \pageref{#1}}
\renewcommand*{\chapterheadstartvskip}{\vspace*{0pt}}					% Abstand vor Kapitelnamen
\renewcommand*{\chapterheadendvskip}{\vspace{1cm}}						% Abstand nach Kapitelnamen
\setlength{\headheight}{30pt}											% Höhe der Kopfzeile festlegen
\setlength{\footskip}{30pt}												% Abstand von Text zu Fusszeile festlegen

% Titelseite
%
% Dokumentation
%

% 1. Seite
\thispagestyle{empty}

\begin{figure}
\flushright						% rechtsbündig
\includegraphics[width=5cm,clip=]{Bilder/HHN_ab_40_mm_4c_pos4}\\
\end{figure}

\vspace*{1.0cm}

\begin{center}
{\Huge Benutzerdokumentation\\
\vspace{0.1cm} 
für die Software\\
\vspace{0.1cm} 
\textbf{we}Factor}

\vspace{3.0cm}

	{\Large Als Projektarbeit für die Vorlesung\\
\vspace{0.1cm}
	Labor für Softwareentwicklung und Project Skills\\
\vspace{0.1cm}
	im Studiengang Software Engineering\\
\vspace{0.1cm}
	der Hochschule Heilbronn\\

}
\end{center}

\vspace*{4.5cm}





Autor:\\
{\Large
Florian Wohlgemuth\\

 }
\vspace{1,5cm}
Heilbronn, \today
%\textbf{FIRMA}, Str. 1, 12345 Ort\\
%Betreuer: Dr. Jemand\\

\fancyhf{} 											% clear all header and footer fields
\renewcommand{\headrulewidth}{1pt}					% Die Dicke der Kopfzeilen-Trennline wird auf 1pt festgelegt
\renewcommand{\footrulewidth}{1pt}					% Die Dicke der Fusszeilen-Trennline wird auf 1pt festgelegt
\fancypagestyle{plain}{}
\fancyhead[R]{}										% Kopfzeile Rechts
\fancyfoot[R]{\thepage}								% Fusszeile Rechts

\pagenumbering{Roman}  								% Die Verzeichnisse werden entsprechend mit römischen Seitenzahlen beschriftet

% +++++++++++++++++++++++++++++++++++++ab hier beginnen (vor dem inhaltsverzeichnis)

% Aenderungshistorie
\begin{tabularx}{\textwidth}{|l|X|l|l|} \hline
       \textbf{Version}  	& \textbf{Name}   	& \textbf{Kapitel}  & \textbf{Datum}\\ \hline
       0.1           		& Florens Hückstädt & Kapitel 1 		& 31.12.20014\\
        \hline
\end{tabularx}

% Inhaltsverzeichnis
%
% Dokumentation
%

% Kapitel dem Inhaltsverzeichnis hinzufügen
\addcontentsline{toc}{chapter}{Inhaltsverzeichnis}			% 1. Argument {toc} = Table Of Contents
															% 2. Argument {chapter} = Ebene bestimmen!
															% 3. Argument {Name} = Name des Eintrags
\noindent

\tableofcontents
	\thispagestyle{plain}
	\newpage				% Neue Seite erzwingen

% Abkürzungsverzeichnis
%
% Dokumentation
%

% Kapitel dem Inhaltsverzeichnis hinzufügen
\addcontentsline{toc}{chapter}{Abkürzungsverzeichnis}			% 1. Argument {toc} = Table Of Contents
																% 2. Argument {chapter} = Ebene bestimmen!
																% 3. Argument {Name} = Name des Eintrags
\noindent

\thispagestyle{plain}

\chapter*{Abkürzungsverzeichnis}
\begin{acronym}[PP]
\acro{JRE}{Java Runtime Environment}



\end{acronym}


% Abbildungsverzeichnis
%
% Dokumentation
%

% Kapitel dem Inhaltsverzeichnis hinzufügen
\addcontentsline{toc}{chapter}{Abbildungsverzeichnis}			% 1. Argument {toc} = Table Of Contents
																% 2. Argument {chapter} = Ebene bestimmen!
																% 3. Argument {Name} = Name des Eintrags
\noindent

\thispagestyle{plain}


\listoffigures



\pagenumbering{arabic}  							% Der folgende Text wird mit arabischen Zahlen beschriftet



% Kapitl Überblick

\thispagestyle{plain}

\chapter{Überschrift}
Lorem ipsum dolor sit amet, consetetur sadipscing elitr, sed diam nonumy eirmod tempor invidunt ut labore et dolore magna aliquyam erat, sed diam voluptua. At vero eos et accusam et justo duo dolores et ea rebum. Stet clita kasd gubergren, no sea takimata sanctus est Lorem ipsum dolor sit amet. Lorem ipsum dolor sit amet, consetetur sadipscing elitr, sed diam nonumy eirmod tempor invidunt ut labore et dolore magna aliquyam erat, sed diam voluptua. At vero eos et accusam et justo duo dolores et ea rebum. Stet clita kasd gubergren, no sea takimata sanctus est Lorem ipsum dolor sit amet.





% Literaturverzeichnis
%
% Dokumentation
%

% Kapitel dem Inhaltsverzeichnis hinzufügen
%\addcontentsline{toc}{chapter}{Literaturverzeichnis}			% 1. Argument {toc} = Table Of Contents
																% 2. Argument {chapter} = Ebene bestimmen!
																% 3. Argument {Name} = Name des Eintrags
\noindent

\thispagestyle{plain}

%\chapter*{Literaturverzeichnis}
\bibliographystyle{apacite}
\bibliography{techDoku}  % Bibtex-Datei wird schon in der Preambel eingebunden





% Anhang
%\include{16_Anhang}

\end{document}